% ============================================================================
% Modelo de ARTIGO CIENTÍFICO — Instituto de Ensino Superior (iCEV)
% ============================================================================
% COMO USAR NO OVERLEAF:
%   1. Faça upload desta pasta inteira (main.tex + referencias.bib) como um
%      novo projeto.
%   2. Edite os campos marcados com "TODO" abaixo (título, autor, curso).
%   3. Escreva o conteúdo de cada seção no lugar dos parágrafos de exemplo.
% ============================================================================

\documentclass[
	article,
	12pt,
	oneside,
	a4paper,
	section=TITLE,
	english,
	brazil,
	sumario=tradicional
	]{abntex2}

\usepackage[T1]{fontenc}
\usepackage[utf8]{inputenc}
\usepackage{indentfirst}
\usepackage{color}
\usepackage{graphicx}
\usepackage{microtype}

\usepackage{mathptmx}

\renewcommand{\ABNTEXsectionfont}{\rmfamily\bfseries}
\renewcommand{\ABNTEXsectionfontsize}{\fontsize{18}{21.6}\selectfont}

\usepackage[brazilian,hyperpageref]{backref}
\usepackage[alf]{abntex2cite}

\renewcommand{\backrefpagesname}{Citado na(s) página(s):~}
\renewcommand{\backref}{}
\renewcommand*{\backrefalt}[4]{
	\ifcase #1 %
		Nenhuma citação no texto.%
	\or
		Citado na página #2.%
	\else
		Citado #1 vezes nas páginas #2.%
	\fi}%

\newcommand{\logoicevheader}{\includegraphics[height=1.9cm]{imagens/logo-icev}}

\makepagestyle{icevpretextual}
\makeoddhead{icevpretextual}{}{\logoicevheader}{}

\makepagestyle{icevheadings}
\makeoddhead{icevheadings}{}{\logoicevheader}{}
\makeoddfoot{icevheadings}{}{\thepage}{}

\setlength{\headheight}{2.0cm}
\setlength{\headsep}{0.5cm}
\checkandfixthelayout

\providecommand{\imprimirsubtitulo}{}
\newcommand{\subtitulo}[1]{\renewcommand{\imprimirsubtitulo}{#1}}

% TODO: preencha com os dados reais do seu trabalho. Em \autor, informe o
% nome do aluno e, na linha seguinte (separados por \\), o nome do professor
% orientador. Se o trabalho não tiver subtítulo, apague a linha
% \subtitulo{} e o \par correspondente dela dentro do \maketitle logo abaixo.
\titulo{Geração de Workflows RPA a partir de Linguagem Natural com LLMs}
\subtitulo{avaliação experimental entre geração restrita e livre por schema em aplicações web}
\autor{Antonio Santana Castelo Branco Bisneto \\ Nome do Professor Orientador Completo}
\local{Teresina -- PI}
\data{08/09/2026}

\renewcommand{\maketitle}{%
	\thispagestyle{icevpretextual}
	\begin{center}
		{\fontsize{18}{21.6}\selectfont\bfseries\MakeUppercase{\imprimirtitulo}\par}
		{\fontsize{18}{21.6}\selectfont\imprimirsubtitulo\par}
	\end{center}
	\vspace{2cm}
	\begin{flushright}
		\fontsize{14}{16.8}\selectfont\imprimirautor
	\end{flushright}
	\vspace{2.5cm}
}

\renewenvironment{resumoumacoluna}[1][\resumoname]{%
	\PRIVATEbookmarkthis{#1}%
	\parindent0pt\relax
	\noindent\textbf{#1:}\par
	\vspace{12pt}
}{%
	\par
	\PRIVATEclearpageifneeded
}

\definecolor{blue}{RGB}{41,5,195}
\makeatletter
\hypersetup{
		pdftitle={\@title},
		pdfauthor={\@author},
		pdfsubject={Artigo científico},
		pdfcreator={LaTeX with abnTeX2},
		pdfkeywords={icev}{artigo}{cientifico},
		colorlinks=true,
		linkcolor=blue,
		citecolor=blue,
		filecolor=magenta,
		urlcolor=blue,
		bookmarksdepth=4
}
\makeatother

\setlength{\parindent}{1.25cm}
\setlength{\parskip}{0pt}

\makeindex

\begin{document}

\selectlanguage{brazil}
\frenchspacing
\pagestyle{icevpretextual}

\maketitle

% Resumo: parágrafo único, de 150 a 500 palavras (NBR 6022:2018).
\begin{resumoumacoluna}
A Automação Robótica de Processos (RPA) tem se consolidado como uma
abordagem relevante para a execução de tarefas repetitivas e baseadas em
regras, mas sua adoção enfrenta desafios recorrentes na etapa de
construção das automações, que tradicionalmente exige converter
manualmente um processo já conhecido em uma sequência precisa de
interações com a interface de um sistema. Este trabalho investiga em que
medida modelos de linguagem de grande escala (LLMs) podem ser utilizados
para gerar workflows RPA executáveis em aplicações web a partir de
requisitos descritos em linguagem natural, propondo uma arquitetura na
qual a saída do LLM é traduzida em uma representação intermediária
estruturada e submetida a um mecanismo determinístico de validação
sintática e semântica antes de ser executada por um motor próprio
construído sobre o Playwright. O eixo experimental do trabalho compara
dois modos de geração a partir da mesma descrição textual: um modo
restrito por schema formal (\textit{constrained}) e um modo sem restrição
estrutural, interpretado de forma tolerante (\textit{free}), ambos
submetidos ao mesmo validador, isolando o efeito específico da restrição
estrutural sobre a saída do LLM. Para avaliar a solução, foi construído um
catálogo de dezoito cenários de teste, distribuídos em três níveis de
complexidade, executados contra uma aplicação-alvo própria e controlada,
totalizando 108 ensaios por rodada de coleta. As métricas analisadas
incluem taxa de validade da geração, taxa de sucesso de execução,
aderência aos requisitos, tempo de geração, tipos de erro e necessidade de
intervenção humana. Os resultados obtidos permitem discutir em que
condições cada modo de geração se mostra mais adequado e em que medida a
combinação entre LLM e validação determinística é suficiente para
viabilizar a geração automatizada de RPAs web.

\vspace{\onelineskip}
\noindent
\textbf{Palavras-chave:} \textit{Automação Robótica de Processos}. \textit{Modelos de Linguagem de Grande Escala}. \textit{Geração de Workflows Web}.
\end{resumoumacoluna}

% Abstract: opcional pela NBR 6022, mas comum no modelo do iCEV.
\renewcommand{\resumoname}{Abstract}
\begin{resumoumacoluna}
\begin{otherlanguage*}{english}
Versão completa do resumo em língua inglesa.

\vspace{\onelineskip}
\noindent
\textbf{Keywords:} \textit{keyword 1}. \textit{keyword 2}. \textit{keyword 3}.
\end{otherlanguage*}
\end{resumoumacoluna}

\textual
\pagestyle{icevheadings}

\section{Introdução}

Robotic Process Automation (RPA) é uma abordagem de automação voltada
principalmente à execução de tarefas repetitivas, estruturadas e baseadas em
regras, permitindo que bots de software reproduzam interações normalmente
realizadas por usuários em sistemas digitais. A literatura aponta RPA como
uma área de crescente relevância para organizações, ao mesmo tempo em que
identifica desafios relacionados à seleção de processos, desenvolvimento,
manutenção, governança e evolução das automações \cite{wewerka2020,
syed2020}. Entre esses desafios, a própria construção inicial da automação
concentra parte relevante do esforço, pois exige converter manualmente um
processo já conhecido em uma sequência precisa de interações com a
interface do sistema-alvo.

Nos últimos anos, o avanço da inteligência artificial generativa e dos
Modelos de Linguagem de Grande Escala (LLMs) ampliou as possibilidades de
interação entre usuários e sistemas de automação. Em vez de exigir que toda
a lógica de um processo seja especificada manualmente, modelos de linguagem
podem interpretar descrições textuais de processos e produzir representações
estruturadas \cite{bettencourt2026}. Essa possibilidade já vem sendo
explorada diretamente no contexto de RPA: \citeonline{chen2026} propõem o
AutoRPA, no qual agentes baseados em LLM sintetizam funções de automação
reutilizáveis a partir da observação de interações com a interface gráfica.
A aplicação industrial reforça a relevância do problema: o Microsoft Power
Automate já disponibiliza recursos de Copilot capazes de receber descrições
em linguagem natural e gerar estruturas de automação
\cite{microsoftcopilot2026}. Esse cenário demonstra que o tema está
alinhado a uma tendência tecnológica real --- mas também que a contribuição
deste trabalho não pode ser apresentada como a simples ideia de gerar
automações por linguagem natural, e sim como a arquitetura adotada, o
recorte de uma representação intermediária estruturada submetida a
validação determinística, e a avaliação sistemática da solução.

O problema de pesquisa que orienta este trabalho é: em que medida modelos de
linguagem de grande escala podem ser utilizados para gerar workflows RPA
executáveis em aplicações web a partir de requisitos descritos em linguagem
natural? Parte-se da hipótese de que a utilização de um LLM existente,
combinada com uma representação intermediária estruturada e um mecanismo
determinístico de validação e execução, pode reduzir o esforço necessário
para construir RPAs web, mantendo uma taxa de sucesso suficiente para
demonstrar a viabilidade da abordagem em processos controlados.

O objetivo geral deste trabalho é desenvolver e avaliar uma plataforma que
utilize LLMs existentes para interpretar requisitos em linguagem natural e
gerar workflows RPA executáveis em aplicações web. De forma específica,
buscou-se: (i) definir um modelo estruturado para representar ações e
etapas de um workflow RPA; (ii) integrar um LLM existente via API para
converter requisitos em linguagem natural em workflows estruturados; (iii)
desenvolver um mecanismo de validação capaz de impedir ou rejeitar ações não
suportadas pela plataforma; (iv) implementar um motor de execução de RPA
para aplicações web utilizando Playwright; (v) registrar informações de
execução (sucesso, falhas, duração, etapas); (vi) construir cenários de
teste com diferentes níveis de complexidade; e (vii) avaliar a solução por
métricas de taxa de sucesso, aderência aos requisitos, tempo de geração,
erros e necessidade de intervenção humana.

Este artigo está organizado em seis seções, incluindo esta introdução. A
Seção 2 apresenta o referencial teórico, abordando os conceitos e os
desafios de RPA, o uso de LLMs na modelagem e na geração de processos e os
trabalhos relacionados à geração de automações a partir de linguagem
natural. A Seção 3 descreve a metodologia, detalhando a arquitetura da
plataforma desenvolvida, os dois modos de geração comparados, o mecanismo
de validação, o motor de execução baseado em Playwright e o ambiente
experimental controlado. A Seção 4 apresenta os resultados obtidos com a
execução dos cenários de teste nos modos \texttt{constrained} e
\texttt{free}. A Seção 5 discute esses resultados à luz do referencial
teórico. Por fim, a Seção 6 apresenta as considerações finais, retomando o
problema de pesquisa e a hipótese que orientaram o trabalho.

\section{Referencial teórico}

Apresente os conceitos e explicações fundamentadas que sustentam a
discussão do tema. É importante que haja uma estrutura lógica, permitindo
ao leitor compreender as ideias convergentes e divergentes entre os
autores trabalhados e o pesquisador.

Citações de até 3 linhas devem estar no corpo do texto, entre aspas; já
citações com mais de 3 linhas devem ser apresentadas em bloco recuado, com
fonte 10 e espaçamento simples, como no ambiente \texttt{citacao} a
seguir:

\begin{citacao}
Citações com mais de três linhas devem ser destacadas com recuo de 4 cm da
margem esquerda, fonte tamanho 10 e espaçamento simples entre linhas,
conforme a NBR 10520 \cite{exemplo2024}.
\end{citacao}

\section{Metodologia}

Trata-se de uma pesquisa aplicada, de natureza experimental, conduzida por
meio do desenvolvimento de uma plataforma e da execução controlada de
cenários de teste. O objetivo não é apenas verificar se a geração de
workflows por linguagem natural é tecnicamente possível --- o que já é
demonstrado por soluções industriais como o Copilot do Power Automate e
por trabalhos como o AutoRPA \cite{microsoftcopilot2026, chen2026} --- mas
avaliar sistematicamente uma arquitetura específica: um formato
intermediário validado deterministicamente entre a saída do LLM e a
execução real, e o efeito de restringir ou não essa saída por um schema
formal.

\subsection{Arquitetura da plataforma}

A plataforma é composta por dois componentes desacoplados, que se
comunicam exclusivamente por meio de uma API HTTP: um front-end para
cadastro de automações e visualização de gerações, execuções e métricas; e
um back-end responsável por três funções centrais --- interpretar a
descrição em linguagem natural por meio de um LLM, validar a saída
produzida e executá-la de fato em um navegador controlado.

Internamente, o back-end é organizado em módulos com responsabilidades
isoladas: um módulo de esquema, responsável apenas pela forma dos dados e
pelas regras de validação; um módulo de planejamento, responsável somente
pela comunicação com o LLM; um motor de execução, responsável por
interpretar um workflow validado e delegar cada ação a um manipulador
específico; e uma camada de abstração sobre o mecanismo de automação de
navegador, que permite substituir a implementação real por uma
implementação simulada em testes. Essa separação permite, por exemplo,
trocar o provedor de LLM sem alterar o motor de execução, e testar o motor
de execução sem depender de um navegador real.

\subsection{Representação intermediária do workflow}

O artefato central produzido pelo LLM não é código executável diretamente,
mas um workflow estruturado, representado como uma lista ordenada de
ações. Foram definidos nove tipos de ação (\texttt{navigate}, \texttt{click},
\texttt{fill}, \texttt{select}, \texttt{waitFor}, \texttt{extract},
\texttt{forEach}, \texttt{if} e \texttt{download}), cada um com um
conjunto fixo de campos permitidos. As ações \texttt{forEach} e \texttt{if}
admitem sub-passos aninhados, com profundidade máxima limitada, para
representar iteração e decisão condicional sem exigir uma estrutura de
controle turing-completa.

Uma decisão arquitetural central é que o LLM nunca é autorizado a gerar
seletores de interface na forma de CSS ou XPath. Todo alvo de interação é
descrito de forma semântica --- por rótulo, papel (\textit{role}), texto
visível, \textit{placeholder} ou identificador de teste ---, o que é
mapeado diretamente para os métodos de localização semântica do
Playwright. Essa escolha tem por objetivo reduzir a fragilidade da
automação frente a mudanças estruturais da interface, e concentrar a
maior parte dos erros de geração em uma única categoria (erro de seleção),
facilitando sua análise.

\subsection{Modos de geração comparados}

O eixo experimental central deste trabalho é a comparação entre dois
modos de geração do workflow a partir da mesma descrição em linguagem
natural:

\begin{description}
\item[Modo \texttt{constrained}] --- a chamada ao LLM é acompanhada de um
schema formal (derivado da própria definição das ações), que restringe
estruturalmente o que o modelo pode produzir como saída.
\item[Modo \texttt{free}] --- a chamada ao LLM não impõe restrição
estrutural; o modelo é instruído, por meio do prompt, a responder apenas
com um objeto JSON, e a saída é interpretada de forma tolerante.
\end{description}

Em ambos os modos, a saída --- estruturalmente restrita ou não --- é
submetida ao mesmo mecanismo de validação, descrito a seguir. Essa
simetria é o que permite isolar, nos resultados, o efeito específico da
restrição estrutural: qualquer diferença observada entre os dois modos não
decorre de regras de validação distintas, mas exclusivamente da presença
ou ausência do schema no momento da geração.

\subsection{Mecanismo de validação}

Todo workflow produzido pelo LLM, em qualquer modo, passa por um
validador determinístico antes de ser considerado apto à execução. A
validação ocorre em duas camadas: uma camada sintática, que verifica se a
estrutura do JSON produzido corresponde ao formato definido para cada tipo
de ação; e uma camada semântica, que aplica regras adicionais não
capturadas apenas pela forma dos dados --- entre elas, a restrição de que
URLs de navegação e valores de seletor pertençam ao domínio autorizado
para aquela automação, a restrição de que valores preenchidos em campos
somente contenham marcadores de credencial de uma lista permitida (nunca
texto livre interpretado como segredo), e o limite máximo de profundidade
de aninhamento para \texttt{forEach} e \texttt{if}.

Um workflow reprovado em qualquer uma das duas camadas é marcado como
inválido e não é executado; os erros de validação, os tokens consumidos e
a latência da chamada ao LLM são registrados mesmo nesse caso, por
constituírem dado experimental relevante para a comparação entre os
modos.

\subsection{Motor de execução}

Workflows validados são executados por um motor próprio, construído sobre
o Playwright, em um navegador \textit{headless} isolado por execução. O
motor percorre os passos do workflow em sequência, interrompendo a
execução no primeiro passo que falhar, e delega cada ação a um manipulador
específico responsável por traduzi-la em uma interação real com a página.
Valores de credencial referenciados por marcadores (como usuário e senha)
são resolvidos a partir de variáveis de ambiente do próprio motor de
execução, nunca sendo expostos ao LLM ou persistidos como texto plano. O
resultado de cada passo --- situação, duração, mensagem de erro quando
houver, evidência em imagem em caso de falha, e valor eventualmente
extraído --- é registrado individualmente, permitindo reconstruir
integralmente o comportamento de uma execução após o fato.

\subsection{Ambiente experimental controlado}

Para garantir a reprodutibilidade dos experimentos, todos os cenários de
teste têm como alvo uma aplicação web própria e isolada, desenvolvida
especificamente para este trabalho, executada localmente e não sujeita a
alterações externas --- ao contrário de um site de terceiros, cuja
interface poderia mudar entre rodadas de coleta e invalidar a comparação.

Foi definido um catálogo fixo de dezoito cenários de teste, distribuídos
igualmente em três níveis de complexidade: simples (fluxo linear de
navegação, preenchimento, extração e exportação), intermediário
(acrescenta filtro ou paginação e uma estrutura de controle) e complexo
(combina múltiplas estruturas de controle e estágios). Um dos cenários do
nível complexo foi deliberadamente definido como não representável pela
plataforma --- uma forma de iteração aninhada que a aplicação-alvo não
expõe --- com o objetivo específico de observar o comportamento de cada
modo de geração diante de um requisito impossível de satisfazer; por essa
razão, esse cenário é analisado separadamente e não é contabilizado no
cálculo da taxa de sucesso geral.

Cada um dos dezoito cenários é executado nos dois modos de geração, com
três repetições por combinação cenário-modo, número escolhido como
compromisso entre mitigar a variabilidade inerente à geração por LLM e
manter viável o custo e o tempo de coleta em um trabalho de conclusão de
curso. Isso totaliza 108 ensaios por rodada completa de coleta
($18 \times 2 \times 3$). Para cada ensaio, são registrados: validade sintática e semântica
do workflow gerado; sucesso ou falha da execução; aderência do workflow
executado aos requisitos do cenário, avaliada por um critério
determinístico que verifica a presença das ações obrigatórias, o
cumprimento de um limite de passos e a correspondência dos valores
extraídos com os valores esperados; tempo de geração; tempo de execução;
tipos de erro observados; e consumo de tokens. A única métrica não
extraída automaticamente é a necessidade de intervenção humana, registrada
manualmente durante a revisão dos casos em que o workflow gerado exigiu
correção para se tornar utilizável.

\section{Resultados}

Apresente os resultados obtidos na pesquisa, utilizando tabelas, gráficos
e figuras quando ajudarem o leitor a compreender os dados encontrados.

\section{Discussão}

Discuta os resultados apresentados, relacionando-os com o referencial
teórico e apontando convergências, divergências e possíveis explicações
para os achados.

\section{Considerações finais}

Nesta seção, faça as considerações finais sobre o trabalho, voltando ao
problema central e aos objetivos apresentados anteriormente. Avalie o que
foi proposto e confirme ou refute as hipóteses levantadas. Por fim, revise
brevemente os conceitos principais discutidos e os autores que
fundamentaram o artigo, sem introduzir novas ideias, autores ou citações.
A conclusão deve ser objetiva, refletindo sobre o que foi abordado ao
longo do trabalho.

\bookmarksetup{startatroot}
\postextual

% TODO: edite o arquivo referencias.bib com as fontes reais do seu trabalho.
\bibliography{referencias}

\end{document}
